\documentclass{article}
\usepackage{graphicx} % Required for inserting images
\usepackage{geometry}
\geometry{a4paper, left=20mm, right =20mm, top=25mm}
\usepackage{parskip} 
\usepackage{transparent} 
\usepackage{mathtools, amssymb, amsfonts, bm, nicefrac}
\usepackage{tabularx, tabularray, booktabs}
\usepackage{float, placeins}
\restylefloat{figure}

\title{\textbf{Methodology (Ongoing)}}
\author{Brayan Mora}
\date{}

\begin{document}

\maketitle
\begin{center}
    \textit{Project Status: Active Research Phase.  While the core statistical foundations and ethical safeguards of this study are firmly established, specific sections, particularly the Section 5: Triangulation Strategy, are undergoing final refinements as we integrate behavioral and geographical overlaps. We invite reviewers to check back weekly for the most current version. Once the ``Ongoing" tag is removed, the methodology will be considered the final Version 1.0.}
\end{center}

\section{Data Source and Harmonization Strategy}
\subsection{Dataset Selection and Scope}
The primary data source for this analysis is the National Household Survey (ENAHO) conducted by the instituto Nacional de Estadistica e Informatica (INEI). Specifically, the analysis utilizes module 500 (Employment and Income) and Module 01 (Household Characteristics)
\begin{itemize}
    \item \textbf{Time Horizon}: Data was pooled from the period 2017-2024. The 2017 cutoff was selected to ensure consistency in survey instrumentation (specifically regarding ethnicity and stratification changes) and to maximize data quality. 
    \item \textbf{Pooling Rationale}: A pooled cross-sectional approach was adopted to increase the effective sample size ($N > 100,000$ observations) This aggregations is critical for the LNOB approach, as identifying ``furthest behind" populations often requires analyzing small subsample (e.g., specific intersections of ethnicity, gender, and labor conditions) that would be statistically invisible in a single-year dataset. 
\end{itemize}
\subsection{Economic Harmonization (Inflation Adjustment) and Main Economic Variable}
The core variable for the analysis is the ``Liquid Income" variable. According to the ENAHO (Empleo e Ingresos) survey, liquid income is the total income (including extra hours, bonuses, Payment for refreshments, transport, commissions, etc.) (source). The liquid income is the total income minus taxes, legal deductions, Pension system(AFP,
ONP, Military and Police Fund), and other deductions (legal, associations, cooperative loans, banks, etc.). This variable can be found as p524e1. However, the variable in used is the deflated-imputed variable. The variable is the i524e1. For deflated only values, the variable is d524e1. 
To validate the use of pooled data, a deflation and harmonization process was implemented to express all monetary values in 2024 Constant Soles.
\begin{enumerate}
    \item \textbf{Intra-Year Deflation}: We utilized the survey-supplied imputed deflated values (i524e1) which adjust for spatial price differences within Peru relative to a base year (2011).
    \item \textbf{Inter-Year Harmonization}: To account for national inflation over the 8-year period, we constructed a Harmonization Factor using the Consumer Price Index (CPI) collected from the ``Banco Central de Reserva del Perú (BCRP)". 
    \begin{itemize}
        \item Calculation: The annual CPI average was calculated for each year ($CPI_y$). The harmonization factor ($H_y$) was derived as:
        \begin{equation*}
            H_y=\frac{CPI_{2024}}{CPI_y}
        \end{equation*}
        \item \textbf{Application}: The final adjusted income ($I_{adj}$) for an individual in year $y$ is calculated as:
        \begin{equation*}
            I_{adj} = I_{survey} \times H_y
        \end{equation*}
        \item This ensures that an income earned in 2017 is comparable in real purchasing power to an income earned in 2024. 
    \end{itemize}
\end{enumerate}
\section{Geospatial Feature Engineering}
To operationalize the ``Geographic Barrier" dimension, a novel variable measuring the distance between residence and workplace was engineered. 
\begin{itemize}
    \item \textbf{Geolocation}: We utilized the UBIGEO coding system (6-digit Department-Province-District codes). Spatial centroid for all districts were extracted using shapefiles from the UN OCHA repository, supplemented with data from the National Geographic Institute (IGN) for missing districts (12060, and 120604).
    \item \textbf{The Haversine Calculation}: The distance was computed using the Harvesine formula, which accounts for the curvature of the earth, calculating the ``as-the-crow-flies" distance between the centroid of the District of Residence and the District of Work:
    \begin{equation*}
        d = 2r\text{ }arcin\left(\sqrt{sin^2\left(\frac{\phi_2-\phi_1}{2}\right) +cos(\phi_1)cos(\phi_2) sin^2\left(\frac{\lambda_2-\lambda_1}{2}\right)}\right)
    \end{equation*}
    Where $\phi$ is latitude, $\lambda$ is longitude, and $r$ is the earth's radius.
\end{itemize}
\section{The Safe-Fail Statistical Framework}
This study adopts a conservative “safe-fail” analytical framework designed to prioritize estimate stability, interpretability, and ethical restraint over maximal disaggregation. The framework is explicitly diagnostic rather than inferential, and its purpose is to prevent false precision, unstable subgroup rankings, and unintended disclosure when analyzing highly intersected populations in complex survey data.
\subsection{Complex Survey Design}
All statistical estimates were calculated using the survey package in R to ensure national representativeness. The design object explicitly accounted for:
\begin{itemize}
    \item Primary Sampling Units (PSU): Clusters (conglome)
    \item Strata: Geographical Layering (estrato)
    \item Weights: Expansion Factors (fac500a)
\end{itemize}
\subsection{The relative Inequality Metric}
Rather than relying solely on absolute income gaps, we utilized a \textbf{Relative Purchasing Power Gap} metric to measure the depth of inequality
\begin{equation*}
    Gap_{ratio}=\frac{Median_{ref}-Median_{target}}{Median_{Ref}}
\end{equation*}
Where $Median_{ref}$ is the median income of the reference group (e.g., Men) and $Median_{target}$ is the median income of the target group (e.g., Women).
\subsection{Reliability Protocols (``Safe-Fail" Filter)}
 To ensure data integrity, every calculated gap was subjected to a rigorous automated filter. An estimate was marked as ``Unreliable" and discarded if it failed any of the following conditions:
\begin{enumerate}
    \item Sample Size: Subgroup $n \geq 100$ per gender
    \item Statistical Significance: The 95\% Confidence interval of the Gap included zero $(p>0.05)$
    \item Coefficient of Variation (CV): 
    \begin{itemize}
        \item Reliable: $CV \leq 15\%$
        \item Acceptable: $15\% < CV \leq 30\%$
        \item Unreliable: $CV > 30\%$
    \end{itemize}
\end{enumerate}
These criteria are applied jointly, as no single metric is sufficient to ensure reliability in highly intersected survey subpopulations; estimates must simultaneously meet minimum size, statistical distinguishability, and variance-stability conditions to be retained. \\
The minimum subgroup size threshold of n $\geq$ 100 per gender functions as a reporting and algorithm-continuation criterion rather than as a claim of sufficient statistical power. In sequential intersectional disaggregation, small cell sizes can generate highly unstable estimates that are disproportionately influenced by single observations, leading to spurious path selection and misleading subgroup rankings. The threshold therefore serves three complementary purposes: (i) to reduce estimate volatility in median-based gap calculations, (ii) to prevent algorithmic over-selection of extreme but substantively fragile intersections, and (iii) to mitigate disclosure and stigmatization risks associated with narrowly defined social groups. The value of 100 reflects a conservative, design-based balance between analytical granularity and ethical restraint rather than a universal statistical requirement.
\section{The Randomized ``Deprived" Algorithm}
\subsection{The intersectional Search Problem}
A central challenge in LNOB analysis is the ``Combinatorial Explosion." With over 7 key structural variables (Ethnicity, Region, Pay Frequency, Contract, Age, Education, Commuting), these are hundreds of possible intersectional subgroups. Manually testing these introduces researcher bias. 
\subsection{Algorithm Solution}
We developed a custom Iterative Random-path Disaggregation Algorithm to objectively identify the groups facing the most severe structural exclusion. 
\begin{enumerate}
    \item Randomized Selection: At each node of the decision tree, the algorithm randomly selects a subset of 5 variables. This prevents the model from deterministically choosing the same variable (e.g., Ethnicity) first every time, ensuring broader exploration of the solution space. 
    \item Recursive Optimization: The algorithm calculates the income gap for all subgroups within the random subset. It selects the ``Deprived" branch (lowest median income + highest inequality) and repeats the process on that subgroup. 
    \item Pathfinding \& Replication: This process is repeated with 500 replications. The final output identifies the ``Most Frequent Deprived Paths" –the specific combinations of identities (e.g., Indigenous $\to$ Daily Pay $\to$ Rural) that consistently appear at the bottom of the socioeconomic ladder across all simulations. 
\end{enumerate}
\section{Triangulation Strategy} (IN REVISION. I have not yet made a decision on the best way to approach how to identify the deprived)
To validate the algorithm findings, the analysis employs a Three-Pillar Triangulation approach: 
\begin{enumerate}
    \item Intersectional Profiling: Deep-dive profiling of the top 3 ``Deprived" groups identifies by the algorithm (analyzing their specific barriers to health, trust, and rights.)
    \item Individual Profiling: Profiling individually the top variables that appears the most in the algorithm. 
    \item Geospatial Bivariate Mapping: A macro-level validation mapping Median Income vs. Inequality at the departmental and natural regional levels to visually explore the ``Zones of Exclusion"
\end{enumerate}
\section{Data Ethics and Privacy}

\end{document}